\section{跨政权回望：制度的收益、代价与史料的沉默}

宋、辽、西夏、金、南宋、蒙古与元并非依次登场的几幅朝代肖像，而是在道路、河流、草场、城镇、市场和海域上不断互相改变的政治体。宋的开封中枢与东南财赋、辽的多中心网络、西夏的绿洲与城寨、金对华北的接收、南宋的江海财政、元的行省与站赤，各自把资源转成行动。它们都必须回答同一个难题：如何让远方的人、粮、钱、信息和军队在不压垮地方承受力的条件下到达需要它们的地方。

\subsection{制度的共同收益}

这些国家留下的并不只是疆界更替。考试、文书、法律、译写、仓储、货币、市场管理、驿传和地方行政，使统治者能够在更大的尺度上识别并协调问题；寺院、宗族、商帮、军户和城市社群则在国家之外或边缘维系生活与互助。多政权竞争有时迫使各方改善运输、边防和外交技术，和平时期的使节、互市和岁币也把冲突限制在可谈判的框架内。征服帝国扩大了人员、技术和商品流动的范围，令许多地区以前所未有的方式相连。

\subsection{制度的共同代价}

同一套能力会制造新的脆弱。把军队和财赋集中于首都，能降低地方军头坐大的风险，也会使道路中断或都城失守格外致命；把赋役、纸币和市场接入国家账目，能扩大调度，却可能把财政紧张转化为地方债务、物价或差役；把边界制度化，能减少全面战争，却要以长期驻军、堡寨和外交成本为代价。战争与征服从不只在战场结束：迁徙、户籍改造、土地重分、城镇毁坏和档案散失，会在数十年间改变家庭能够选择的生计。

因此，不能把“统一”当作自动的制度进步，也不能把“分裂”当作单纯的失败。辽、宋与西夏的并立曾维持长期而昂贵的和平；金与南宋在战争和和议之间形成新的区域秩序；蒙古、元的扩张提供广阔的联系，同时以沉重征调与不均衡接收为条件。每种秩序都有受益者、承担者与无法进入官书的人。

\subsection{怎样阅读本卷的证据}

本卷的事件锚点保留在各政权文件中，供索引和查证使用；它们不再承担连续阅读的负担。读者在正文看到的是由事件簇组织的选择、后果与比较，而不是把每个年号条目重新排印一遍。编年材料仍不可替代：它们使陈桥、澶渊、永乐城、靖康、绍兴、襄阳、崖山和元末河工能够置于相对可靠的时间关系中。可是日期可靠，并不等于军数、动机、社会影响和每句传闻都已可靠。

《宋史》《辽史》《金史》《元史》提供王朝叙事的骨架，却多为后出的官修史；《续资治通鉴长编》《建炎以来系年要录》《宋会要辑稿》让宋代政事与制度更为稠密，也各有编纂、残缺和立场。黑水城文书、契约、碑刻、墓志、考古、钱币、地方志及多语材料能让地方、宗教、市场和迁徙显影，却通常只保存某一地点和一部分人。材料不一致时，本卷不以最动人的故事填满空白，而是说明争点来自何处，并把可证、可推与不可断言分开。

下一卷进入明代时，不能把元末只当作一个应当被清除的前史。行省、江南财赋、边疆军政、海上贸易、社会流动和灾害应对所遗留的问题，会在新的王朝中以不同方式继续出现。历史的连续性不在于一项制度原封不动地存留，而在于后来的行动者仍要面对它曾经解决、也曾经制造的限制。
